\QHeading{第1题　脱式计算，能简算的要简算。}
\par\noindent


\CalcBlock[方法一]{1}{\begin{aligned}[t]&\phantom{={}}\color{black}7.45-1.31-3.69\\&\color{answerred}=7.45-(1.31+3.69)\\&\color{answerred}=7.45-5\\&\color{answerred}=2.45\end{aligned}}
\par
\MethodBlock{方法二：按运算顺序直接计算（对照）}{本法只作结果核对；本题要求简便计算，前面的简便解法仍保留。}{\begin{aligned}[t]&\phantom{={}}7.45-1.31-3.69\\
&=6.14-3.69\\
&=2.45\end{aligned}}


\CalcBlock[方法一]{2}{\begin{aligned}[t]&\phantom{={}}\color{black}6.3+9.85+12.7\\&\color{answerred}=6.3+12.7+9.85\\&\color{answerred}=(6.3+12.7)+9.85\\&\color{answerred}=19+9.85\\&\color{answerred}=28.85\end{aligned}}
\MethodBlock{方法二：先处理括号或乘除，再凑整再减去多加的数}{把第二个数凑成整数，再减去多加的部分。}{\begin{aligned}[t]&\phantom{={}}6.3+9.85+12.7\\
&=16.15+12.7\\
&=16.15+13-0.3\\
&=29.15-0.3\\
&=28.85\end{aligned}}

\CalcBlock[方法一]{3}{\begin{aligned}[t]&\phantom{={}}\color{black}9.56-(3.5+3.56)\\&\color{answerred}=9.56-(3.56+3.5)\\&\color{answerred}=9.56-3.56-3.5\\&\color{answerred}=6-3.5\\&\color{answerred}=2.5\end{aligned}}
\MethodBlock{方法二：去括号后计算}{加号后的括号去掉不变号，减号后的括号去掉要改变各项的加减号。}{\begin{aligned}[t]&\phantom{={}}9.56-(3.5+3.56)\\
&=9.56-3.5-3.56\\
&=6.06-3.56\\
&=2.5\end{aligned}}

\MethodBlock{方法三：按运算顺序直接计算（对照）}{本法只作结果核对；本题要求简便计算，前面的简便解法仍保留。}{\begin{aligned}[t]&\phantom{={}}9.56-(3.5+3.56)\\
&=9.56-7.06\\
&=2.5\end{aligned}}

\CalcBlock[方法一]{4}{\begin{aligned}[t]&\phantom{={}}\color{black}7.32-(2.68+2.32)\\&\color{answerred}=7.32-5\\&\color{answerred}=2.32\end{aligned}}
\MethodBlock{方法二：去括号后计算}{加号后的括号去掉不变号，减号后的括号去掉要改变各项的加减号。}{\begin{aligned}[t]&\phantom{={}}7.32-(2.68+2.32)\\
&=7.32-2.68-2.32\\
&=4.64-2.32\\
&=2.32\end{aligned}}

\CalcBlock{5}{\begin{aligned}[t]&\phantom{={}}\color{black}8.92-0.08+5.72\\&\color{answerred}=8.84+5.72\\&\color{answerred}=14.56\end{aligned}}
\CalcBlock{6}{\begin{aligned}[t]&\phantom{={}}\color{black}12.85-5.15+7.62\\&\color{answerred}=7.70+7.62\\&\color{answerred}=15.32\end{aligned}}
\CalcBlock[方法一]{7}{\begin{aligned}[t]&\phantom{={}}\color{black}83.67-(51.3-7.5)\\&\color{answerred}=83.67-43.8\\&\color{answerred}=39.87\end{aligned}}
\MethodBlock{方法二：去括号后计算}{加号后的括号去掉不变号，减号后的括号去掉要改变各项的加减号。}{\begin{aligned}[t]&\phantom{={}}83.67-(51.3-7.5)\\
&=83.67-51.3+7.5\\
&=32.37+7.5\\
&=39.87\end{aligned}}

\CalcBlock{8}{\begin{aligned}[t]&\phantom{={}}\color{black}14.08+5.29-6.44\\&\color{answerred}=19.37-6.44\\&\color{answerred}=12.93\end{aligned}}
\CalcBlock{9}{\begin{aligned}[t]&\phantom{={}}\color{black}17.5+2.56-4.47\\&\color{answerred}=20.06-4.47\\&\color{answerred}=15.59\end{aligned}}
\par\noindent

\QHeading{第2题　解决问题。}

\Problem{① 一根绳子长15米，第一次用去6.25米，第二次用去5.75米，还剩下多少米？}
\par


\MethodTitle{方法一}
\Answer{{\Formula{\begin{aligned}[t]&\phantom{={}}15-6.25-5.75\\&=15-(6.25+5.75)\\&=15-12\\&=3\text{（米）}\end{aligned}}}\par 答：还剩3米。}

\MethodTitle{方法二：按先后顺序连续减}
\Answer{先求第一次用后剩下多少，再减第二次用的。}
\Answer{\Formula{\begin{aligned}[t]&\phantom{={}}15-6.25-5.75\\
&=8.75-5.75\\
&=3\text{（米）}\end{aligned}}}
\Answer{答：还剩3米。}
\Problem{② 小兰原有9.2元，卖废品挣了3.8元，买练习本花了3.2元。小兰现在有多少钱？}

\MethodTitle{方法一：简便计算}
\Answer{{\Formula{\begin{aligned}[t]&\phantom{={}}9.2+3.8-3.2\\&=(9.2-3.2)+3.8\\&=6+3.8\\&=9.8\text{（元）}\end{aligned}}}\par 答：小兰现在有9.8元。}

\MethodTitle{方法二：先求这次收支净增加的钱}
\Answer{卖废品挣的钱减去买练习本的钱，再加原来的钱。}
\Answer{\Formula{\begin{aligned}[t]&\phantom{={}}9.2+(3.8-3.2)\\
&=9.2+0.6\\
&=9.8\text{（元）}\end{aligned}}}
\Answer{答：小兰现在有9.8元。}
